\documentclass[../../main.tex]{subfiles}

\begin{document}

\section{$C_{c}(X)$ 上的正线性泛函与 Riesz-Markov 表示定理}\label{section:ch21-rieszmarkov}

令 $X$ 为局部紧 Hausdorff 空间而 $C_{c}(X)$ 为 $X$ 上具有紧支撑的连续实值函数的线性空间. 线性泛函 $\Lambda:C_{c}(X)\to\mathbf{R}$ 称为是正的, 若
\begin{align*}
\text{对 }C_{c}(X)\text{ 中的所有 }f,\quad f\geqslant 0\Rightarrow\Lambda(f)\geqslant 0.
\end{align*}

\begin{theorem}[Riesz-Markov 表示定理]\label{theorem:ch21-rieszmarkov}
令 $X$ 为局部紧 Hausdorff 空间而 $\Lambda$ 是 $C_{c}(X)$ 上的正线性泛函. 则存在 $\mathcal{B}(X)$ 上的 Radon 测度 $\mu$ 使得
\begin{align}\label{align:ch21-rieszmarkov等式}
\text{对所有 }f\in C_{c}(X),\quad\Lambda(f)=\int_{X}f\,d\mu.
\end{align}
此外, 可以选取 $\mu$ 使得它是内正则的, 即对每个开集 $\mathcal{O}$,
\begin{align*}
\mu(\mathcal{O})=\sup\{\mu(K)\mid K\subseteq\mathcal{O},\ K\text{ 是紧的}\}.
\end{align*}
在 Radon 测度中, $\mu$ 是唯一的.
\end{theorem}

\begin{proof}
我们首先假设 $X$ 是紧的. 定义集函数 $\lambda$ 在开集族 $\mathcal{T}$ 上为
\begin{align*}
\text{对 }\mathcal{O}\in\mathcal{T},\quad\lambda(\mathcal{O})=\sup\{\Lambda(f)\mid f\in C_{c}(X),\ 0\leqslant f\leqslant 1,\ \operatorname{supp}f\subseteq\mathcal{O}\}.
\end{align*}
我们验证 $\lambda$ 满足 Radon 测度构造定理的假设.
\begin{enumerate}[(i)]
\item 单调性显然.
\item 次可加性: 令 $\{\mathcal{O}_{k}\}_{k=1}^{\infty}$ 为开集族且令 $\mathcal{O}=\bigcup_{k=1}^{\infty}\mathcal{O}_{k}$. 取 $f\in C_{c}(X)$, $0\leqslant f\leqslant 1$, $\operatorname{supp}f\subseteq\mathcal{O}$. 由于 $\operatorname{supp}f$ 是紧的, 存在有限 $n$ 使得 $\operatorname{supp}f\subseteq\bigcup_{k=1}^{n}\mathcal{O}_{k}$. 令 $\{\varphi_{k}\}_{k=1}^{n}$ 为从属于 $\{\mathcal{O}_{k}\}_{k=1}^{n}$ 的单位分解. 则 $f=\sum_{k=1}^{n}\varphi_{k}f$, 且 $\Lambda(f)=\sum_{k=1}^{n}\Lambda(\varphi_{k}f)\leqslant\sum_{k=1}^{n}\lambda(\mathcal{O}_{k})\leqslant\sum_{k=1}^{\infty}\lambda(\mathcal{O}_{k})$. 取上确界得到 $\lambda(\mathcal{O})\leqslant\sum_{k=1}^{\infty}\lambda(\mathcal{O}_{k})$.
\item 不交开集的可加性: 令 $\mathcal{O}_{1}$ 和 $\mathcal{O}_{2}$ 为具有不相交闭包的开集. 令 $f_{1}$ 和 $f_{2}$ 为 $C_{c}(X)$ 中的函数, $0\leqslant f_{i}\leqslant 1$, $\operatorname{supp}f_{i}\subseteq\mathcal{O}_{i}$. 则 $f_{1}+f_{2}$ 的支撑包含于 $\mathcal{O}_{1}\cup\mathcal{O}_{2}$, 且 $\Lambda(f_{1}+f_{2})=\Lambda(f_{1})+\Lambda(f_{2})$. 因此 $\lambda(\mathcal{O}_{1}\cup\mathcal{O}_{2})\geqslant\lambda(\mathcal{O}_{1})+\lambda(\mathcal{O}_{2})$. 反过来, 若 $f$ 的支撑包含于 $\mathcal{O}_{1}\cup\mathcal{O}_{2}$, 我们可以用 Urysohn 引理构造 $g_{1}$ 和 $g_{2}$ 使得 $g_{i}\in C_{c}(X)$, $0\leqslant g_{i}\leqslant 1$, $\operatorname{supp}g_{i}\subseteq\mathcal{O}_{i}$, 且 $g_{1}+g_{2}=1$ 在 $\operatorname{supp}f$ 上. 则 $f=f g_{1}+f g_{2}$, 且 $\Lambda(f)=\Lambda(f g_{1})+\Lambda(f g_{2})\leqslant\lambda(\mathcal{O}_{1})+\lambda(\mathcal{O}_{2})$. 因此 $\lambda(\mathcal{O}_{1}\cup\mathcal{O}_{2})\leqslant\lambda(\mathcal{O}_{1})+\lambda(\mathcal{O}_{2})$.
\end{enumerate}
根据 Radon 测度构造定理, 存在 $\mathcal{B}(X)$ 上的 Radon 测度 $\mu$ 使得对每个开集 $\mathcal{O}$, $\mu(\mathcal{O})=\lambda(\mathcal{O})$. 我们验证\eqref{align:ch21-rieszmarkov等式}.

令 $f\in C_{c}(X)$. 我们证明 $\Lambda(f)=\int f\,d\mu$. 由于 $\Lambda$ 和积分都是线性的, 我们可以假设 $0\leqslant f\leqslant 1$. 取 $\varepsilon>0$ 且令 $n$ 为自然数. 对 $k=0,1,\cdots,n$, 定义
\begin{align*}
\mathcal{O}_{k}=\{x\in X\mid f(x)>k/n\}.
\end{align*}
则 $\mathcal{O}_{0}=X$ 且 $\mathcal{O}_{n}=\varnothing$. 令 $K_{k}=\{x\in X\mid f(x)\geqslant k/n\}$. 选取 $g_{k}\in C_{c}(X)$ 使得 $0\leqslant g_{k}\leqslant 1$, $g_{k}=1$ 在 $K_{k}$ 上, 且 $\operatorname{supp}g_{k}\subseteq\mathcal{O}_{k-1}$. 定义 $f_{k}=g_{k}-g_{k+1}$ (其中 $g_{0}=1$, $g_{n+1}=0$). 则 $\operatorname{supp}f_{k}\subseteq\mathcal{O}_{k-1}$ 且 $f_{k}\geqslant 0$. 我们有
\begin{align*}
f=\sum_{k=1}^{n}(k/n)f_{k}+\rho,
\end{align*}
其中 $0\leqslant\rho\leqslant 1/n$. 因此
\begin{align*}
\Lambda(f)=\sum_{k=1}^{n}(k/n)\Lambda(f_{k})+\Lambda(\rho)\leqslant\sum_{k=1}^{n}(k/n)\mu(\mathcal{O}_{k-1})+\varepsilon.
\end{align*}
另一方面,
\begin{align*}
\int f\,d\mu\geqslant\sum_{k=1}^{n}(k/n)\mu(K_{k})\geqslant\sum_{k=1}^{n}(k/n)[\mu(\mathcal{O}_{k})-\varepsilon].
\end{align*}
取 $n$ 充分大得到 $\Lambda(f)\leqslant\int f\,d\mu+\varepsilon$. 类似的论证给出反向不等式. 因此 $\Lambda(f)=\int f\,d\mu$.

对于一般的局部紧 Hausdorff 空间 $X$, 我们考虑它的 Alexandroff 单点紧化 $X^{*}$. 对 $f\in C_{c}(X)$, 我们可以将它延拓为 $X^{*}$ 上的连续函数 $\widehat{f}$ 使得 $\widehat{f}(\omega)=0$. 定义 $\widehat{\Lambda}(\widehat{f})=\Lambda(f)$. 则 $\widehat{\Lambda}$ 是 $C(X^{*})$ 上的正线性泛函. 根据紧情形的证明, 存在 $\mathcal{B}(X^{*})$ 上的 Radon 测度 $\widehat{\mu}$ 使得 $\widehat{\Lambda}(\widehat{f})=\int_{X^{*}}\widehat{f}\,d\widehat{\mu}$. 定义 $\mu(E)=\widehat{\mu}(E)$ 对 $E\in\mathcal{B}(X)$. 则 $\mu$ 是 $\mathcal{B}(X)$ 上的 Radon 测度且\eqref{align:ch21-rieszmarkov等式}成立.
\end{proof}

\subsection*{习题}

\begin{exercise}\label{exercise:ch21-28}
令 $X$ 为紧 Hausdorff 空间而 $\Lambda$ 是 $C(X)$ 上的正线性泛函. 证明存在常数 $M$ 使得对所有 $f\in C(X)$, $|\Lambda(f)|\leqslant M\cdot\|f\|_{\max}$.
\end{exercise}

\begin{exercise}\label{exercise:ch21-29}
在 Riesz-Markov 定理的证明中, 验证 $\lambda$ 的次可加性.
\end{exercise}

\begin{exercise}\label{exercise:ch21-30}
令 $X$ 为局部紧 Hausdorff 空间而 $\mu$ 是 $\mathcal{B}(X)$ 上的 Radon 测度. 证明由 $\Lambda(f)=\int f\,d\mu$ 定义的泛函 $\Lambda$ 是 $C_{c}(X)$ 上的正线性泛函.
\end{exercise}

\end{document}